% !TEX root = ../main_v2.tex
\section{Discussion}
\label{sec:discussion}

This paper has three distinct but interlocking contributions: a structural
criterion for operator representability of the HMF, an exact closed-form
solution in the transverse-field qubit model, and a numerical characterisation
of the crossover between the weak- and ultrastrong-coupling regimes.
We discuss each in turn before drawing out the broader implications.

\subsection*{Closure and representability}

The central algebraic result of Sec.~\ref{sec:closure} is that the question
``does a local closed-form HMF exist?'' has a precise answer in terms of the
operator algebra generated by repeated adjoint action of $H_Q$ on the coupling
$f$.
Conditions (C1)--(C3) distil this into a hierarchy of closure requirements:
adjoint closure of the coupling chain $\{f_n\}$, closure of that chain's
quadratic span (which contains the Gaussian influence operator $\Delta$),
and BCH closure of commutators among the quadratic elements.
When all three hold, $H_{\mathrm{MF}}(\beta)$ is a finite operator polynomial
and can be extracted by a single BCH
logarithm~\cite{magnusExponentialSolutionDifferential1954a,blanesMagnusExpansionIts2009,weiLieAlgebraicSolution1963,vanbruntVisser2015BCHClosedForms};
when they fail, the
HMF necessarily has support outside the target operator family and any
representation within that family is a controlled approximation.

A notable consequence is the impossibility result for non-quadratic
continuous-variable systems.
For a Hamiltonian $H_Q = p^2/2m + V(x)$ with a polynomial potential
$\deg V > 2$, the adjoint chain $\{f_n\}$ does not close in any
finite-dimensional algebra: the chain grows without bound and generates
the full universal enveloping algebra.
This obstruction is analogous in spirit to the Groenewold--Van Hove theorem
on consistent polynomial quantisation, and extends it to the equilibrium
context: for such potentials no finite local operator can
represent the exact reduced equilibrium state.

On the constructive side, the BCH form
Eq.~\eqref{eq:HMF_BCH_def_final} provides a systematic generator of the
HMF in any regime.
The separation of operator structure (carried entirely by the adjoint chain)
from all bath and temperature dependence (entering only through the scalar
kernel moments $C^\gtrsim_{nm}(\beta)$) means that, in practice, one need
only compute these moments --- themselves one-dimensional Laplace-type
integrals over the spectral density --- and then perform matrix algebra
in the operator space.
This is computationally tractable even when the adjoint chain is long, and
suggests a natural numerical scheme: truncate the chain at a finite depth,
compute the corresponding moments, and use the BCH formula to assemble the
approximate HMF.
The approximation is then purely on the operator side; all bath physics
enters exactly.

\subsection*{The exact qubit solution}

The spin-boson model is the canonical testbed for open-quantum-system physics,
yet the exact HMF in the transverse-coupling case has not, to our knowledge,
been written down explicitly outside of weak-coupling or ultrastrong limiting
forms~\cite{cresserWeakUltrastrongCoupling2021a,hovhannisyanHamiltonianMeanForce2020,millerHamiltonianMeanForce2018}.
The result~\eqref{eq:HMF_final_clean} achieves this using only three bath
numbers: the kernel functional $\mathcal{K}$ evaluated at zero frequency,
at the system resonance, and in the resonant channel $\mathcal{R}(\omega_q)$.

The compression of the entire bath into three scalars is exact,
not perturbative, and follows directly from the $\mathfrak{su}(2)$ closure
of the adjoint chain.
The algebra guarantees that no higher-frequency evaluations can appear:
all imaginary-time convolutions collapse, via the ordered Green's kernel
$\mathcal{G}^>(x,y)$ and its closed-form reduction~\eqref{eq:Gxy_closed},
to a finite set of Laplace transforms evaluated on the system's own mass shell.
This is the imaginary-time analogue of the fact that, in linear response, a
two-level system couples to the bath only at its transition frequency.

The crossover parameter $\chi = \sqrt{\Delta_z^2 + \Sigma^2}$ that emerges
from the $\mathfrak{su}(2)$ algebra is more than a convenient shorthand.
Its square is the total accumulated bath response seen by the qubit in
imaginary time --- the analogue of the polaron displacement energy ---
and $\gamma(\chi) = \tanh\chi/\chi$ is the universal renormalisation factor
that interpolates between the two fixed points.
At $\chi \ll 1$, $\gamma \approx 1$ and the state is a weakly dressed
bare Gibbs distribution.
At $\chi \gg 1$, $\gamma \sim 1/\chi$ and the state is pinned to the
interaction direction, reproducing the ultrastrong fixed point of
Cresser and Anders~\cite{cresserWeakUltrastrongCoupling2021a,hovhannisyanHamiltonianMeanForce2020} exactly.
The crossover between these regimes is smooth, parameter-free, and
analytically located at $\chi = 1$ --- the precise condition $g = g_\star(\beta)$.

This solution also sheds light on an apparent paradox that arises in the
intermediate steps of the calculation.
The influence operator $\Delta(\beta)$ has manifestly asymmetric ladder
amplitudes, $\Sigma_+ \neq \Sigma_-$, which would seem to imply a
non-Hermitian reduced state.
The resolution is that the asymmetry is a gauge artefact of computing
the influence before recombining with the bare Gibbs factor $e^{-\beta H_Q}$:
the KMS relation Eq.~\eqref{eq:Sigma_pm_kms_relation_v5} precisely maps
$\Sigma_+ = e^{\beta\omega_q}\Sigma_-$, and the product $e^{-\beta H_Q}e^\Delta$
is Hermitian by construction.
The FDT is not an additional input but a consistency condition that
is automatically enforced by the imaginary-time kernel.

\subsection*{The crossover as a representability bottleneck}

The numerical results of Sec.~\ref{sec:numerical_results} add a
dimension to the physical picture that is not visible from the analytic
solution alone: the sensitivity of the simulated state to the 
representational limits of the Hilbert space.

The distinction between the discrete and continuous analytic branches
in Fig.~\ref{fig:ed_vs_analytic} is physically sharp.
Both branches evaluate the same exact formula; crucially, the 
discrete-analytic branch (which uses the same mode set as ED) correctly 
tracks the continuum limit as $\chi \to \infty$. This proves that the 
HMF theory is exact for any mode set, and that a discrete bath possesses 
the same intrinsic ability to reach the ultrastrong limit as the continuum.

The divergence of ED from its own discrete analytic prediction as 
$\chi \gtrsim 1$ therefore reveals a **representability bottleneck**.
In the strongly coupled regime, the bath oscillators acquire large 
polaron-like displacements. Representing such highly occupied states in 
a simulation with a fixed Fock cutoff $n_{\max}$ is fundamentally impossible; 
the truncated Hilbert space cannot accommodate the tails of the displaced 
ground state, forcing the simulation to stay at an artificially high energy 
level. This is a "mode occupation" failure, not a discretization failure.

The bandwidth analysis of Fig.~\ref{fig:bandwidth_convergence} makes this 
quantitative. Choice of bandwidth $B$ determines how the displacement "load" 
needed to reach the US limit is distributed among the available modes. 
The sharp peak in cutoff sensitivity is a direct marker of the available 
modes hitting their Fock boundaries as they attempt to polarize. This is 
not an arbitrary convergence artifact; it is the physical signature of 
trying to represent a highly non-local equilibrium state (in terms of mode 
occupancy) within a restricted representational space. The crossover 
parameter $\chi$ governs both the analytic structure of the state and the 
boundary where this representability bottleneck becomes insurmountable.

\subsection*{Connections to the literature}

This work sits at the intersection of several lines of inquiry.

The Cresser-Anders framework~\cite{cresserWeakUltrastrongCoupling2021a,hovhannisyanHamiltonianMeanForce2020}
gives the leading-order weak-coupling and exact ultrastrong-coupling forms
for the qubit mean-force state.
The present result interpolates between these limits exactly, identifies
$\chi$ as the relevant dimensionless variable throughout, and gives the
crossover scale $g_\star(\beta)$ in closed form.

The Du-Chen finite-bath HMF~\cite{duGeneralizedHamiltonianMeanforce2025a}
generalises the mean-force concept to finite reservoirs by introducing
a pair of effective Hamiltonians for system and bath.
Our discrete-analytic branch is the commuting Gaussian instance of this
construction: replacing the continuum by a discrete mode set is precisely
the finite-bath limit, and the parametric sensitivity we observe in
Fig.~\ref{fig:bandwidth_convergence} is a concrete numerical realisation
of the bath-feedback effects they discuss analytically.

The Burke-Iles-Smith structural
study~\cite{burkeStructureHamiltonianMean2024} analyses the operator
content of the HMF through correlation-function methods.
The potential-renormalisation
perspective~\cite{correaPotentialRenormalisationLamb2025} provides a complementary
bath-induced operator characterisation.
Our closure criterion (C1)--(C3) provides a further algebraic
characterisation: operator growth in the HMF corresponds precisely to
the failure of one or more closure conditions.

More broadly, methods that access $\bar\rho_Q(\beta)$ numerically ---
HEOM, stochastic Schrödinger/Liouville
approaches~\cite{tanimuraReducedHierarchicalEquations2014,stockburgerSimulatingSpinbosonDynamics2004},
chain mappings, and reaction-coordinate schemes~\cite{nazirReactionCoordinateMapping2018,IlesSmithDijkstraLambertNazir2016,AntoSztrikacsNazirSegal2023,ShubrookIlesSmithNazir2025} --- are all,
in a sense, performing the quenched-density average numerically.
For a recent unifying overview connecting these approaches, see
Ref.~\cite{XuVadimovStockburgerAnkerhold2026}.
Our BCH/adjoint-chain organisation shows how to extract a closed-form
operator representation from such an average whenever the algebra
closes, and identifies the minimal set of bath evaluations required
(three numbers in the qubit case).
This suggests a hybrid programme: use numerical methods to compute kernel
moments on a target spectral density, then assemble the HMF analytically
via the BCH formula.
Surveys of the growing thermodynamic strong-coupling programme can be found
in Refs.~\cite{campbellRoadmapQuantumThermodynamics2026,binderThermodynamicsQuantumRegime2018}.

\subsection*{Limitations and outlook}

The central analytical restriction of the present work is Gaussianity.
The quenched representation and the cumulant expansion hold for any bath,
but the exact collapse of $\Delta$ to a bilinear form in the adjoint chain
relies on the bath cumulant hierarchy truncating at second order.
Beyond the Gaussian sector, higher cumulants $K^{(n)}$ with $n \geq 3$
enter~\eqref{eq:influence_cumulant_expansion} and introduce products
$f_{n_1}\cdots f_{n_k}$ of arbitrarily high order.
Closure of such a family in a finite algebra is generically much harder.
The non-Gaussian extension therefore represents the natural next step,
and the adjoint-chain organisation already provides the skeleton: each
additional cumulant order contributes a new layer of kernel moments
multiplying operator monomials of corresponding degree, and the BCH
structure persists.

Within the Gaussian sector, the qubit model studied here is the
simplest closed algebra.
The next natural target is multi-qubit systems, where the adjoint chain
generated by $H_Q$ will generically involve $k$-body Pauli strings.
Closure then requires the target family to contain those strings, and the
kernel-moment computation becomes a matrix problem in an enlarged operator
space.
Whether the HMF remains a two-body operator for a chain of qubits, or
generates genuinely $k$-local terms, is a question our framework can
in principle answer.

Finally, the crossover physics identified here has immediate consequences
for the growing experimental literature on strong-coupling thermodynamics
with superconducting circuits and semiconductor quantum dots~\cite{campbellRoadmapQuantumThermodynamics2026,binderThermodynamicsQuantumRegime2018}.
The crossover scale $g_\star(\beta) = \chi_0(\beta)^{-1/2}$ translates
directly into a device-parameter threshold: below it, the bare Gibbs
description is adequate; above it, the mean-force correction is
non-perturbative.
The quantitative agreement between the exact formula and ED simulation in
the weak-coupling sector provides a parameter-free benchmark that can be
used to validate strong-coupling simulation protocols in those settings.


